\songtitle{Demasiado rola}

\setin{2005-08-22}
\tag*{2026}\tag{bambuco}\tag{pasillo}\tag{spanish-lyrics}\tag{ricochets}\tag{fabian-muller}\tag{carolina-muller}
% -> ricochets.tex
\begin{notes}
Lyrics by Claude based on \emph{Ricochets} by Carlos Thompson -- chapter 12 (\emph{Vertrek}): in Fabián Müller's voice, the last breakfast before Carolina leaves, changua against calentado montañero (``Me saliste demasiado rola''), the quiet late morning over the crossword, and her question in German answered in the newspaper margin. Retrieved from rataflechera's Suno page (V6, 20 November 2026).
\end{notes}
\begin{style}
Colombian Andean bambuco, tiple and bandola picking over fingerpicked nylon guitar, light bass and soft shaker; warm and playful in D major, around 100 BPM in lilting 6/8, slowing into a minor pasillo bridge in B minor with muted tiple and lone bandola; warm, wry male baritone with a paisa accent, male tenor color, intimate, half-smiling legato croon in carefully Antioqueño Spanish and German phrases; warm close-miked strings, a touch of slapback echo on the voice, final refrain darker and hushed, then a solo tiple fade.
\end{style}
\begin{lyrics}
\songpart{Intro}
\annotation{tiple and bandola, cups and plates in a kitchen, a newspaper unfolding}

\songpart{Verso 1}
Golpea Alma despacito en la puerta del cuarto,
"Señorita Carolina, ya está listo el desayuno."\\
Sale descalza, la toalla en el pelo,
"¿Qué me tienes, Almita?" y se le ilumina el mundo:\\
changüita con almojábanas. ¡Qué rico!, dice ella.\\
Yo la miro desde mi plato sin decir ni una palabra:\\
calentado montañero, un chorizo, un tinto,\\
lo que manda la montaña.

\songpart{Verso 2}
Leche hervida con un huevo ahogado,\\
cebolla larga y cilantro nadando,\\
el calado remojándose, blando y pasmado,\\
como un náufrago en la taza flotando.\\
Eso no lo come un paisa ni enfermo,\\
eso se le da al que no tiene dientes.\\
Y mi niña, que estudia en Alemania,\\
se lo toma feliz y caliente.

\songpart{Coro}
Me saliste demasiado rola,\\
le digo como quien se queja de nada,\\
y ella sonríe y, sin pedir permiso,\\
me roba un pedazo de chorizo del plato.\\
Me saliste demasiado rola,\\
de Bogotá hasta el último bocado,\\
y yo, paisa hasta el chorizo,\\
contento de que me lo haya robado.

\songpart{Verso 3}
Se queda la mañana sin forma ninguna,\\
cada uno en lo suyo, en la misma cocina.\\
Yo con el periódico y el esfero en la mano,\\
el crucigrama a medias, ella con su vida.
"Partida, seis letras." Y ella, sin mirarme:
"Salida." Y la escribo, y no pesa nada.
"¿Ya tienes todo empacado o te llevo temprano\\
donde tu mamá?" "Casi todo. Me tienes todo el día."

\songpart{Coro}
Me saliste demasiado rola,\\
le digo como quien se queja de nada,\\
y ella sonríe y, sin pedir permiso,\\
me roba un pedazo de chorizo del plato.\\
Me saliste demasiado rola,\\
de Bogotá hasta el último bocado,\\
y yo, paisa hasta el chorizo,\\
contento de que me lo haya robado.

\songpart{Puente}
\annotation{pasillo lento en menor, bandola sola}
"Papa…", me dice en alemán, y el aire se detiene,\\
casi nunca me habla en alemán, y nunca de esta forma.\\
Primero me pregunta por Sophie, si es en serio,\\
que ojalá la haga feliz, que ojalá yo sea feliz.\\
Y después, más bajito, como quien no quiere saber:
"Stimmt's eigentlich, dass du in Drogengeschäften mitmischst?"\\
Miro las paredes, las ventanas, los rincones callados,\\
y el esfero del crucigrama se me va solo al margen.

\songpart{Coro final}
\annotation{más despacio, oscuro}
"Nein. Das stimmt nicht", le digo en voz alta,\\
y le paso el periódico con la letra apretada:
"So einfach ist das nicht.\\
Es könnten sich Mikrofone im Haus befinden."\\
Me saliste demasiado rola, mija,\\
y yo te salí demasiado callado.\\
Ella lee, levanta los ojos,\\
y el chorizo se queda en el plato.

\songpart{Outro}
\annotation{tiple solo, bambuco que se apaga}\\
Partida, seis letras.\\
Salida.\\
Me tienes todo el día…
\annotation{End}
\end{lyrics}

